\def\lecturenumber{01}
\documentclass[11pt,a4paper]{article}
\usepackage[T1]{fontenc}
\usepackage[utf8]{inputenc}
\usepackage{lmodern}
\usepackage[margin=19mm,top=21mm,bottom=21mm,headheight=15pt]{geometry}
\usepackage{amsmath,amssymb,booktabs,array,graphicx,microtype}
\usepackage[dvipsnames]{xcolor}
\usepackage{enumitem,fancyhdr,listings,tcolorbox}
\usepackage[colorlinks=true,linkcolor=MidnightBlue,urlcolor=MidnightBlue]{hyperref}
\definecolor{courseblue}{HTML}{164B69}
\definecolor{pale}{HTML}{F0F5F7}
\pagestyle{fancy}
\fancyhf{}
\fancyhead[L]{\small\sffamily\color{courseblue}VALIDATED COMPUTING}
\fancyhead[R]{\small\sffamily Lecture \lecturenumber}
\fancyfoot[L]{\footnotesize Expanded lecture notes \textbullet\ October 2026}
\fancyfoot[R]{\footnotesize\thepage\ / 6}
\renewcommand{\headrulewidth}{0.4pt}
\setlength{\parindent}{0pt}
\setlength{\parskip}{5.5pt plus 1pt}
\setlength{\emergencystretch}{1.5em}
\setlist[itemize]{leftmargin=1.5em,itemsep=3pt,topsep=3pt}
\setlist[enumerate]{leftmargin=1.7em,itemsep=5pt,topsep=4pt}
\lstset{language=Python,basicstyle=\small\ttfamily,keywordstyle=\color{courseblue}\bfseries,
 commentstyle=\color{OliveGreen},stringstyle=\color{BrickRed},showstringspaces=false,
 columns=fullflexible,keepspaces=true,frame=single,rulecolor=\color{black!18},
 backgroundcolor=\color{black!2},aboveskip=6pt,belowskip=6pt,breaklines=true,
 xleftmargin=5pt,xrightmargin=5pt,framesep=5pt}
\newcommand{\R}{\mathbb R}
\newcommand{\fl}{\operatorname{fl}}
\newcommand{\abs}[1]{\left|#1\right|}
\newcommand{\heading}[2]{\section*{\color{courseblue}#1}\vspace{-4pt}{\small\sffamily\color{black!65}#2}\par\vspace{3pt}}
\newcommand{\opening}[2]{{\Large\sffamily\bfseries\color{courseblue}Lecture \lecturenumber\quad #1\par}\vspace{5pt}{\small #2}\par}
\newtcolorbox{question}[1]{colback=pale,colframe=courseblue!55,boxrule=0.4pt,
 arc=1mm,left=7pt,right=7pt,top=4pt,bottom=4pt,title={#1},fonttitle=\bfseries\small,
 coltitle=courseblue,colbacktitle=pale}
\newcommand{\reading}[1]{{\footnotesize\textbf{Reading and notebook.} #1\par}}
\hypersetup{pdfauthor={Validated Computing course},pdfsubject={Expanded introductory lecture notes with questions and exercises}}

\hypersetup{pdftitle={Lecture 01: From numerical answers to certificates}}
\begin{document}
\opening{From answers to certificates}{Prerequisites: Python, inequalities and elementary calculus.\newline
Goal: turn a plausible numerical answer into a precise claim with checkable evidence.}
\heading{1. What does a computed number tell us?}{0--12 minutes: predictions, discussion, and vocabulary}
A program returns a finite object: a bit pattern, a printed decimal, a list of points.
The mathematical question usually concerns an exact real number or every point in
an infinite set. Validated computing supplies a justified connection between those
objects. Before discussing precision, ask: \emph{what exact statement are we trying to establish?}

Predict these results individually, then compare explanations with a neighbour.
\begin{lstlisting}
print(0.1 + 0.2 == 0.3)
large = 1e16
print((large + 1.0) - large)
x = 1.4142
print(x*x - 2.0)
\end{lstlisting}
On the binary64 Python environment used in this course, the first result is
\texttt{False} and the second is \texttt{0.0}. The third is close to zero.
The first two expose representation or rounding; the third raises a different
question: how does a small function value constrain the distance to a root?

\textbf{Four levels of statement.}
An \emph{approximation} proposes a value, such as $1.4142$.
An \emph{error estimate} predicts its accuracy under stated assumptions, which
may still need checking. An \emph{enclosure} is a set proved to contain the intended
exact answer. A \emph{certificate} combines a precise claim with finite evidence
and the mathematical argument explaining why the evidence suffices.

For example, ``two runs print the same digits'' is evidence about two computations.
``The positive root of $x^2-2$ belongs to $[1.4142,1.4143]$'' is a claim about an
exact real number. Agreement can suggest the claim, but cannot replace its proof:
two runs may share the same mistake or discard the same information.

\begin{question}{Discuss: which missing assumption matters?}
A plot crosses the horizontal axis; two endpoint evaluations have opposite signs;
an iteration has stopped changing; two precisions agree. Which of these establishes
a root? For each, name one additional condition or piece of evidence you would need.
\end{question}
\textbf{Working rule.} Write the claim before writing the checker. ``Exactly one root
in this interval'' requires existence \emph{and} uniqueness. ``No root here'' requires
exclusion on the entire interval. ``I could not decide'' is a useful outcome, not a
licence to replace a missing proof with a favourable plot.

\textbf{Two jobs for a numerical program.} An exploratory calculation may locate
a promising interval quickly, using plots, ordinary floating-point arithmetic,
or a heuristic search. A validation then checks a sufficient mathematical condition
on that interval. The search need not be trusted if the final checker independently
verifies all required evidence. This separation lets us retain fast exploratory
methods while making the final conclusion depend on a much smaller argument.

\newpage
\heading{2. Identify the input and the error}{12--25 minutes: exact arithmetic and a short input audit}
There are several places for a numerical answer to depart from the intended problem.
Keeping them separate helps us choose the right remedy.
\begin{center}\small
\begin{tabular}{@{}p{0.23\linewidth}p{0.69\linewidth}@{}}\toprule
Source & Example and question to ask \\ \midrule
Input representation & Does the binary float labelled \texttt{0.1} equal the rational $1/10$?\\
Arithmetic rounding & Was an increment discarded during an addition?\\
Truncation & How large is the omitted tail of an infinite sum?\\
Discretization & What can happen between the sampled grid points?\\
Input uncertainty & Is a measurement an exact value or a whole admissible range?\\
Model discrepancy & Does the mathematical model describe the physical system adequately?\\ \bottomrule
\end{tabular}
\end{center}
A proof about a mathematical model does not automatically bound its discrepancy
from the physical world. Increasing arithmetic precision addresses neither an
unmodelled force nor an uncertain sensor reading.

\textbf{Exact audit of a familiar decimal.} The stored binary64 value of \texttt{0.1} is
\[
 q=\frac{3602879701896397}{36028797018963968},\qquad
 q-\frac1{10}=\frac1{180143985094819840}>0.
\]
Both identities concern rational numbers, so Python's integer and rational arithmetic
can check them without using an approximate decimal reference.
\begin{lstlisting}
from fractions import Fraction
intended = Fraction(1, 10)
stored = Fraction.from_float(0.1)
assert stored - intended == Fraction(1, 180143985094819840)
assert Fraction("0.1") == intended
assert Fraction(0.1) == stored
\end{lstlisting}
The last line does not recover the user's intended decimal: the float conversion
has already happened. Also distinguish $1/10$ from a measurement reported as
$0.1\pm0.001$. The latter means an input set, even if its endpoints are represented exactly.

\textbf{Absolute and relative error.} If the intended answer is $y$ and the computed
answer is $\widehat y$, absolute error is $|\widehat y-y|$. Relative error is
$|\widehat y-y|/|y|$, defined only when $y\ne0$. Neither can normally be computed
exactly unless we already know enough about $y$. A proved upper bound $e$ gives
\[
 |\widehat y-y|\le e\quad\Longleftrightarrow\quad
 y\in[\widehat y-e,\widehat y+e].
\]
The endpoints must themselves be computed exactly or rounded outwards.
\begin{question}{Pair question}
An instrument reports a length as $2.000\pm0.005$ metres. Would 100-digit arithmetic
make the length known to 100 digits? Which error source would remain?
\end{question}

\textbf{Choose an error scale for the question.} An absolute error of one unit
is small relatively when the true answer is $10^9$, but enormous when it is
$10^{-9}$. Near a zero, relative error can be misleading or undefined, while an
absolute tolerance remains meaningful. A sensible report states both the error
metric and its units. ``Accurate to six digits'' is incomplete unless the intended
meaning, such as a relative bound or correct decimal rounding, is specified.

\newpage
\heading{3. Build a complete root certificate}{25--45 minutes: board proof followed by an exact checker}
Let $f(x)=x^2-2$. We want to prove that there is exactly one positive root in
\[
 I=[a,b],\qquad a=\frac{14142}{10000},\quad b=\frac{14143}{10000}.
\]
A complete argument has three parts. First, a polynomial is continuous on $I$.
Second, exact integer arithmetic gives
\[
 14142^2=199996164<200000000<200024449=14143^2,
\]
so $f(a)<0<f(b)$. The intermediate value theorem therefore gives a root in $(a,b)$.
Third, $f'(x)=2x>0$ on $I$, so $f$ is strictly increasing and the root there is unique.
The sign checks provide evidence; continuity and monotonicity explain its force.

\begin{lstlisting}
from fractions import Fraction

a = Fraction(14142, 10000)
b = Fraction(14143, 10000)
assert 0 < a < b
assert a*a < 2 < b*b
midpoint = (a + b) / 2
radius = (b - a) / 2
print(midpoint, radius)  # 5657/4000 and 1/20000
\end{lstlisting}
The midpoint is $1.41425$ and the radius is $0.00005$. Consequently the exact
positive root $r$ satisfies $|r-1.41425|\le0.00005$. This enclosure is already a
validated result, although no interval-arithmetic library was needed.

\textbf{Improve the certificate without knowing the root.} Put $m=(a+b)/2$.
If $m^2<2$, replace $a$ by $m$; if $m^2>2$, replace $b$ by $m$. If $m^2=2$,
the midpoint is an exact root and we can stop. The sign bracket is preserved,
and each nonterminal step halves its width. After $n$ steps the width is
$(b_0-a_0)2^{-n}$. The midpoint error is at most half that width. Thus the desired
accuracy determines a finite number of steps in advance.

\begin{question}{Board activity: refine and interpret}
Perform the first two bisections using fractions. Keep a table of the endpoints,
midpoint sign, and interval width. Which statement is the loop invariant?
Does the original bracket alone certify that rounding the root to four decimal
places gives $1.4142$? Locate the rounding threshold before answering.
\end{question}
\textbf{Certificate record.} Save the function, exact endpoints, endpoint signs,
existence theorem, uniqueness argument, and final width. A second reader should
be able to check this record without trusting the code that guessed the bracket.
Our reasoning trusts elementary real analysis and the exact-arithmetic implementation;
it is not a formal verification of the Python interpreter.

\textbf{When may we print certified digits?} To certify rounding to $p$ decimal
places, it is sufficient for the entire enclosure to lie strictly inside one
rounding bin
$((k-\tfrac12)10^{-p},(k+\tfrac12)10^{-p})$ for an integer $k$.
Then every possible exact answer rounds to $k10^{-p}$. A small interval that
straddles a rounding boundary does not decide the rounded answer. Width alone
therefore does not determine which digits are certified.

\newpage
\heading{4. Why common stopping tests are incomplete}{45--60 minutes: two counterexamples and two useful bounds}
\textbf{Residual versus error.} A residual measures how closely an equation is
satisfied. For $h(x)=10^{-12}(x-1)$, the residual at $x=0$ has magnitude $10^{-12}$,
but the distance to the unique root is $1$. Multiplying an equation by a small
constant makes residuals small without moving any root.

A theorem can connect a residual to an error. Suppose $r$ is a known root,
$f$ is continuous on the segment between $r$ and $x$ and differentiable inside it,
and $|f'(t)|\ge m>0$ throughout that segment. The mean value theorem gives
\[
 f(x)-f(r)=f'(\xi)(x-r),\qquad
 |x-r|\le\frac{|f(x)|}{m}.
\]
This is conditional on existence of $r$ and the derivative bound. Neither condition
follows merely from observing a small residual. If the residual is computed with
rounding, we also need a justified upper bound on its exact magnitude.

For the root certified on page 3, $r,x\in[1,2]$ implies $f'(t)\ge2$. With
$x=1.4142$ interpreted as an exact decimal, $|f(x)|=0.00003836$, so
$|x-r|\le0.00001918$. This improves the one-sided endpoint distance bound obtained
from the original bracket. The root-existence proof was an essential earlier step.

\textbf{A positive grid does not prove positivity.} The function
$g(x)=(x-1/3)^2$ vanishes at $1/3$, yet is strictly positive at each point
$0,0.1,\ldots,1$. Even exact evaluation at every point of that grid misses the zero.
The issue is coverage between points, not floating-point accuracy.

Here is one way to bridge that gap. Assume $|g'(x)|\le L$ on $[a,b]$, and let a
uniform grid include both endpoints with spacing $h$. Every point is within $h/2$
of a grid point $x_j$. The mean value theorem yields
\[
 g(x)\ge g(x_j)-L|x-x_j|\ge \min_j g(x_j)-Lh/2.
\]
If the final lower bound is positive, then $g$ is positive everywhere on $[a,b]$.
Use certified lower bounds for the grid values if their evaluations are rounded.

\begin{question}{Apply the theorem, not just the picture}
For $g(x)=x^2+1$ on $[-1,1]$, use $L=2$ and the grid
$-1,-1/2,0,1/2,1$. What global lower bound follows? Is it sharp?
If a bound is not positive, have we proved that a zero exists?
\end{question}

\textbf{Do not omit continuity.} The expression $f(x)=1/x$ has opposite signs
at $-1$ and $1$, but no real zero. It is not defined at the interior point zero,
so the intermediate value theorem does not apply on $[-1,1]$. Sampling more points
without checking the domain does not repair that argument. Root validation must
check the hypotheses of the theorem as carefully as the numerical inequalities.

\newpage
\heading{5. Bound what was never computed}{60--75 minutes: a tail estimate and a combined error budget}
An exact finite computation can still answer the wrong question if it omits an
uncontrolled infinite tail. Consider
\[
 S=\sum_{k=1}^{\infty}\frac1{k^2},\qquad
 S_N=\sum_{k=1}^{N}\frac1{k^2},\qquad R_N=S-S_N.
\]
We do not need a closed form for $S$ to enclose it. Since $x\mapsto x^{-2}$ is
positive and decreasing, comparison with rectangles under and over the curve gives
\[
 \int_{N+1}^{\infty}\frac{dx}{x^2}\le R_N\le
 \int_N^{\infty}\frac{dx}{x^2}.
\]
For the lower inequality, compare $1/k^2$ with the integral over $[k,k+1]$;
for the upper one, compare it with the integral over $[k-1,k]$. Summing from
$k=N+1$ and evaluating the integrals proves
\[
 \boxed{\quad S_N+\frac1{N+1}\le S\le S_N+\frac1N.\quad}
\]
Every endpoint is rational when $S_N$ is accumulated exactly. Notice that the
\emph{enclosure width} is $1/[N(N+1)]$, while the omitted tail itself is of order
$1/N$. Bounding the tail from both sides is more informative than discarding it.
\begin{lstlisting}
from fractions import Fraction

N = 10
partial = Fraction(0)
for k in range(1, N + 1):
    partial += Fraction(1, k*k)
lower = partial + Fraction(1, N + 1)
upper = partial + Fraction(1, N)
assert upper - lower == Fraction(1, N*(N + 1))
print(lower, upper)
\end{lstlisting}
A float sum $\widehat S_N$ would need a separate arithmetic bound. If
$|\widehat S_N-S_N|\le\rho$, then a valid enclosure is
\[
 \left[\widehat S_N-\rho+\frac1{N+1},\;
       \widehat S_N+\rho+\frac1N\right],
\]
provided its endpoints are formed with exact arithmetic or outward rounding.
Truncation control does not secretly include summation error.

\begin{question}{Explain the error budget}
One student increases $N$ but keeps an unreliable summation method. Another computes
$S_{10}$ to 200 digits but ignores the tail. What remains unjustified in each case?
How large must $N$ be for the exact enclosure width to be below $10^{-4}$?
\end{question}

\textbf{Why not use fractions for everything?} Exact rationals give transparent
reference calculations, but their numerators and denominators can grow large.
They also do not directly represent arbitrary square roots or transcendental
values. Later lectures replace exact intermediate values with rigorous enclosures
whose endpoints have controlled finite precision. This trades exact equality for
a proved inclusion, retaining the mathematical guarantee while making more complex
computations practical.

\newpage
\heading{6. Exercises and a certificate audit}{75--90 minutes: choose two problems in pairs; finish the others after class}
For each answer, distinguish observations, hypotheses, and conclusions. A useful
submission includes the intended exact inputs and explains why the checks imply
the stated claim. Code alone is not an existence or uniqueness argument.
\begin{enumerate}
\item \textbf{A new root certificate.} Prove that the positive square root of $3$
lies in $[1732/1000,1733/1000]$. State the theorem for existence and the argument
for uniqueness. Give a midpoint and certified error bound. How many bisections
make the width at most $10^{-6}$? Use integer comparisons, not a call to
\texttt{math.sqrt}, as the evidence.
\item \textbf{Same endpoint signs.} On $[-1,1]$, give three continuous functions
that are positive at both endpoints and have, respectively, no root, exactly one
root, and two distinct roots in the interior. What stronger information would
allow you to exclude a root? Explain why making a plot denser does not supply it.
\item \textbf{Repair a numerical report.} A program states: ``The smallest of
101 sampled values is $0.03$, so $g(x)>0$ on $[0,1]$.'' Assume the uniform grid
includes both endpoints and each value is exact. If $|g'|\le4$, repair the
argument and give a certified lower bound. Repeat with $|g'|\le10$ and report
honestly what the bound does, and does not, decide.
\item \textbf{Combine two bounds.} Suppose $N=100$, a finite-sum approximation is
$\widehat S_N$, and $|\widehat S_N-S_N|\le10^{-8}$. Write a rigorous enclosure
for the infinite sum of page 5. What is its width? Identify the arithmetic that
still needs care when representing its endpoints on a machine.
\end{enumerate}
\begin{question}{Audit another pair's solution}
Can you find the exact claim, domain, and input interpretation? Which facts are
checked by arithmetic and which are supplied by a theorem? Are all theorem
hypotheses stated? What would the program report if one check failed?
\end{question}
\textbf{Selected checks, after attempting the problems.}
For Exercise 1, $1732^2=2999824$ and $1733^2=3003289$;
the midpoint is $1.7325$ and the radius $0.0005$. Ten bisections suffice.
For Exercise 3, the covering radius is $1/200$, giving lower bounds $0.01$
and $-0.02$. The latter is inconclusive, not evidence of a negative value.
For Exercise 4 the width is $1/(100\cdot101)+2\cdot10^{-8}$.

\textbf{Exit question.} Describe a small computation that you could audit without
rerunning the search that produced its answer. What would you save as its certificate?

\reading{Companion: \texttt{notebooks/01\_intro.ipynb}; Lab 01: numerical crime scene.
Tucker, \emph{Validated Numerics} (2011), Chapter 1, especially \S1.2 and \S1.6;
Example 1.6.3 motivates combining a finite sum with a tail bound. The arguments
and exercises above are expanded course material.}
\end{document}
